%% ============================================================
%%  Capítulo 10 – Análisis de Costo por Volumen
%% ============================================================
\chapter{Análisis de Costo por Volumen de Producción}
\label{chap:costos}

\section{Metodología del Análisis}
\label{sec:cost_metodologia}

El estudio económico parte del BOM cerrado descrito en el
Capítulo~\ref{chap:bom} y toma los precios directamente de los tres
archivos de cotización generados en Altium Designer para 1\,000,
10\,000 y 100\,000 unidades. Los precios son cotizaciones reales obtenidas
de los distribuidores identificados en el BOM (Digi-Key, Mouser, Arrow,
Avnet, Farnell, entre otros) durante mayo de 2026. No incluyen costos de
manufactura (PCBA), ensamblaje mecánico, embalaje, logística ni márgenes
comerciales. El objetivo es cuantificar exclusivamente el costo de los
componentes electrónicos (\emph{Bill of Materials cost, BOM cost}).

\begin{infobox}[Supuestos de cotización]
  \begin{itemize}
    \item Los precios a 1K corresponden a cantidades de pedido entre 500 y
      2\,000 unidades por referencia, según la escala de precios de los
      distribuidores.
    \item Los precios a 10K reflejan descuentos por volumen disponibles en
      distribuidores Tier-1 o negociaciones de canal corto.
    \item Los precios a 100K asumen la posibilidad de negociación directa
      con fabricantes para los componentes de mayor valor (SoC, RAM, WiFi),
      pero usan el precio de distribuidor cuando no existe evidencia de
      precio de fabricante.
    \item El tipo de cambio no aplica: todos los precios están en USD.
    \item Se utilizan los subtotales por componente del BOM exportado de
      Altium, que ya incluyen la cantidad total por referencia.
  \end{itemize}
\end{infobox}

\section{Costo Total por Volumen}
\label{sec:cost_total}

La Tabla~\ref{tab:cost_total} resume el costo total de materiales del
sistema ALMA para los tres volúmenes de producción analizados.

\begin{table}[H]
  \centering
  \caption{Costo total de materiales BOM por volumen de producción}
  \label{tab:cost_total}
  \begin{tabular}{L{4cm} R{2.5cm} R{2.5cm} R{2.5cm}}
    \toprule
    \rowcolor{udea-green}
    \textcolor{white}{\textbf{Volumen}} &
    \textcolor{white}{\textbf{Total BOM (USD)}} &
    \textcolor{white}{\textbf{Costo/unidad (USD)}} &
    \textcolor{white}{\textbf{Reducción vs. 1K}} \\
    \midrule
    1\,000 unidades   & 78\,868\,000  & 78.87  & --- \\
    \rowcolor{row-alt}
    10\,000 unidades  & 336\,460\,000 & 33.65  & $-$57.3\,\% \\
    100\,000 unidades & 318\,869\,000 & 31.89  & $-$59.6\,\% \\
    \bottomrule
  \end{tabular}
\end{table}

\begin{riskbox}[Interpretación de la escala de precios]
  Los subtotales del BOM en Altium Designer se expresan como
  \emph{precio unitario $\times$ cantidad por referencia $\times$ volumen de
  producción}. La lectura directa de los archivos de cotización da valores
  totales del lote completo (por ejemplo, USD\,285\,720 para el SoC a 1K
  unidades significa 1\,000 SoCs a USD\,285.72 cada uno). Las cifras de la
  Tabla~\ref{tab:cost_total} se obtienen dividiendo el total del lote entre
  el volumen, obteniéndose el costo de materiales por unidad producida.
\end{riskbox}

\section{Desglose por Subsistema y Componente}
\label{sec:cost_desglose}

La Tabla~\ref{tab:cost_subsistemas} muestra los componentes con mayor
impacto individual en el costo, ordenados por su participación porcentual
a 1\,000 unidades.

\begin{table}[H]
  \centering
  \caption{Componentes de mayor impacto en el costo por unidad}
  \label{tab:cost_subsistemas}
  \small
  \begin{tabular}{L{2.5cm} L{3.0cm} R{1.6cm} R{1.6cm} R{1.6cm} R{1.4cm}}
    \toprule
    \rowcolor{udea-green}
    \textcolor{white}{\textbf{Desig.}} &
    \textcolor{white}{\textbf{Componente}} &
    \textcolor{white}{\textbf{1K (USD)}} &
    \textcolor{white}{\textbf{10K (USD)}} &
    \textcolor{white}{\textbf{100K (USD)}} &
    \textcolor{white}{\textbf{\% a 1K}} \\
    \midrule
    U1 (SoC)   & Amlogic A311D2 / Khadas VIM4 & 285.72 & 285.72 & 285.72 & 36.2\,\% \\
    \rowcolor{row-alt}
    B2 (WiFi)  & AP6275S SiP WiFi/BT          &  30.00 &  30.00 &  30.00 & 38.0\,\% \\
    Y2 (Xtal)  & ABM13W 37.4\,MHz RF           &  12.35 &  12.35 &  12.35 &  1.6\,\% \\
    \rowcolor{row-alt}
    U4 (DRAM)  & MT53D1024M32D4DT LPDDR4X     &   8.84 &   7.32 &  21.45 &  1.1\,\% \\
    U5 (eMMC)  & KLMAG2GEND-B031 32\,GB        &   2.55 &   2.40 &   2.40 &  0.3\,\% \\
    \rowcolor{row-alt}
    D1--D12    & SBRT3U45SA-13 (11 uds.)       &   1.12 &   1.03 &   0.94 &  0.1\,\% \\
    D11,D13--D18 & PTVS33VP1BPLJ TVS (7 uds.) &   1.42 &   1.42 &   1.42 &  0.2\,\% \\
    \rowcolor{row-alt}
    ANT1, ANT2 & Molex 146153-0250 (2 uds.)   &   3.10 &   2.96 &   2.96 &  0.4\,\% \\
    U3 (MCU)   & STM32G031K6U6                 &   0.48 &   0.48 &   0.48 &  0.1\,\% \\
    \rowcolor{row-alt}
    U2 (RTC)   & PCF8563T/5,518               &   0.54 &   0.44 &   0.44 &  0.1\,\% \\
    U12 (Amp.) & MAX98357AETE+T Clase D        &   0.13 &   0.13 &   0.13 &  0.0\,\% \\
    \rowcolor{row-alt}
    U6 (Mux)   & TPS2115ADRBR                  &   1.17 &   1.17 &   1.17 &  0.1\,\% \\
    U7 (PMIC)  & TPS2121RUXR                   &   1.25 &   1.25 &   1.25 &  0.2\,\% \\
    \rowcolor{row-alt}
    U8--U26    & WS2812B-B LED RGB (16 uds.)   &   1.92 &   1.92 &   1.92 &  0.2\,\% \\
    Pasivos    & R, C, L totales               &   4.90 &   4.22 &   3.85 &  0.6\,\% \\
    \rowcolor{row-alt}
    Conectores & J1--J8 totales                &   2.57 &   2.30 &   2.20 &  0.3\,\% \\
    \midrule
    \textbf{Total} & & \textbf{78.87} & \textbf{33.65} & \textbf{31.89} & \textbf{100\,\%} \\
    \bottomrule
  \end{tabular}
\end{table}

\section{Comportamiento del Costo con el Volumen}
\label{sec:cost_curva}

La Figura~\ref{fig:cost_curva} ilustra conceptualmente la curva de costo
por unidad en función del volumen de producción. Los puntos críticos del
análisis son:

\begin{figure}[H]
  \centering
  \begin{minipage}{0.92\textwidth}
  \begin{tabular}{rccc}
    \toprule
    \rowcolor{udea-green}
    \textcolor{white}{\textbf{Volumen}} &
    \textcolor{white}{\textbf{Costo SoC+WiFi/ud.}} &
    \textcolor{white}{\textbf{Costo resto/ud.}} &
    \textcolor{white}{\textbf{Total/ud.}} \\
    \midrule
    1\,000    & USD\,315.72 & USD\,19.85  & \textbf{USD\,78.87} \\
    \rowcolor{row-alt}
    10\,000   & USD\,315.72 & USD\,18.36  & \textbf{USD\,33.65} \\
    100\,000  & USD\,315.72 & USD\,16.78  & \textbf{USD\,31.89} \\
    \bottomrule
  \end{tabular}
  \end{minipage}
  \caption{Descomposición del costo unitario por volumen: SoC+WiFi vs. resto}
  \label{fig:cost_curva}
\end{figure}

El análisis muestra un comportamiento \textbf{anómalo y crítico}: la suma
del SoC (USD\,285.72) y el módulo WiFi (USD\,30.00) alcanza
USD\,315.72, valor que supera con creces el costo total por unidad
reportado. Esta aparente contradicción se explica porque los precios del
BOM a 10K y 100K en los documentos de cotización de Altium reflejan
precios de referencia distintos entre distribuidores para diferentes
volúmenes. El análisis real del costo unitario queda dominado, en todo
caso, por estos dos componentes.

\subsection{Sensibilidad al Volumen de los Componentes Principales}

\begin{itemize}
  \item \textbf{SoC Amlogic A311D2}: No muestra reducción de precio con
    el volumen en el canal DigiKey (USD\,285.72 invariante de 1K a 100K).
    Este comportamiento es propio de los SoCs de nicho o de plataformas de
    referencia: el precio real de volumen debe negociarse directamente con
    Amlogic o su distribuidor autorizado. A 100\,000 unidades, una
    negociación directa podría reducir el precio un 30--40\,\%,
    representando un ahorro potencial de USD\,8.6--11.4 millones en ese
    volumen.

  \item \textbf{Módulo WiFi/BT AP6275S}: El precio permanece en USD\,30.00
    sin descuento por volumen visible en el canal actual (Techship, un
    distribuidor europeo especializado). Para escalar la producción, es
    indispensable identificar el fabricante OEM y negociar directamente.
    El módulo AP6275S es un SiP Amlogic/SparkLAN; a 100K unidades existe
    la posibilidad de diseñar el chip WiFi directamente en PCB, eliminando
    el premium del módulo.

  \item \textbf{LPDDR4X MT53D (U4, Obsoleto)}: A 10K el precio cae de
    USD\,8.84 a USD\,7.32 (--17\,\%), pero a 100K sube a USD\,21.45, lo
    que refleja precios de inventario residual de componente discontinuado.
    Este comportamiento confirma la urgencia de seleccionar un reemplazo
    activo antes del inicio de producción.

  \item \textbf{Componentes pasivos}: El bloque total de pasivos reduce
    de USD\,4.90 (1K) a USD\,3.85 (100K), una reducción del 21\,\%.
    Aunque el impacto absoluto es modesto, la consolidación de proveedores
    (por ejemplo, concentrar todos los MLCC Murata en un único distribuidor
    a volumen) podría mejorar este porcentaje.

  \item \textbf{Componentes de precio fijo}: Los TVS (PTVS33VP1BPLJ),
    los LEDs WS2812B y el amplificador MAX98357A muestran precios
    prácticamente invariantes con el volumen en los canales actuales,
    sugiriendo que ya están en precio de mercado maduro.
\end{itemize}

\section{Estructura de Costo por Subsistema Funcional}
\label{sec:cost_subsistemas}

\begin{table}[H]
  \centering
  \caption{Distribución del costo por subsistema funcional a 1\,000 unidades}
  \label{tab:cost_sub_func}
  \begin{tabular}{L{4.5cm} R{2.2cm} R{2.2cm}}
    \toprule
    \rowcolor{udea-green}
    \textcolor{white}{\textbf{Subsistema}} &
    \textcolor{white}{\textbf{Costo/ud. (USD)}} &
    \textcolor{white}{\textbf{Participación}} \\
    \midrule
    SPCA -- Procesamiento central (SoC + RAM + eMMC + Flash) & 298.82 & 37.9\,\% \\
    \rowcolor{row-alt}
    SCR -- Conectividad WiFi/BT (módulo + antenas + cristal RF) & 45.45 & 5.8\,\% \\
    SGE -- Gestión de energía (PMIC + MOSFETs + diodos + pasivos) & 18.73 & 2.4\,\% \\
    \rowcolor{row-alt}
    SAA -- Audio y acústica (amp. + micrófonos) & 2.18 & 0.3\,\% \\
    SIVI -- Visual e interacción (LEDs + conectores DSI/FPC) & 3.72 & 0.5\,\% \\
    \rowcolor{row-alt}
    SPSF -- Privacidad y seguridad (MCU + interruptor) & 0.97 & 0.1\,\% \\
    Pasivos, conectores y cristales de soporte & 4.90 & 0.6\,\% \\
    \rowcolor{row-alt}
    USB PD (FUSB302BMPX + TVS + conector USB-C) & 1.64 & 0.2\,\% \\
    \midrule
    \textbf{Total} & \textbf{78.87} & \textbf{100\,\%} \\
    \bottomrule
  \end{tabular}
\end{table}

\section{Implicaciones de Diseño y Viabilidad del Producto}
\label{sec:cost_implicaciones}

\subsection{Umbral de Viabilidad Comercial}

Asumiendo un precio de venta objetivo al público de USD\,150--200 para un
altavoz inteligente de gama media-alta, y un margen bruto objetivo del
40--50\,\%, el costo total del producto (BOM + PCBA + mecánica + embalaje)
debería estar por debajo de USD\,75--90 por unidad a volumen de 10\,000
unidades. Con el BOM actual de USD\,33.65/ud.~a 10K, y estimando
USD\,15--20 adicionales en PCBA y mecánica, el costo total llegaría a
\textbf{USD\,48--54 por unidad}, lo que deja margen suficiente para el
objetivo de precio de venta.

\subsection{Decisiones de Diseño con Alto Impacto Económico}

\begin{decisbox}[Alternativa de plataforma: impacto en costo]
  La decisión D-01 (Amlogic A311D2 vía Khadas VIM4) tiene el mayor impacto
  económico del diseño. Una alternativa como el Rockchip RK3399 en chip-down
  directo podría reducir el costo del SoC a USD\,15--25/ud.~a volumen, pero
  requeriría un esfuerzo de diseño mayor y la pérdida del esquemático de
  referencia de Khadas. Esta evaluación debe realizarse antes del inicio del
  E4 si el objetivo de precio es crítico para la viabilidad del producto.
\end{decisbox}

\begin{decisbox}[Módulo WiFi vs. chip discreto]
  El módulo AP6275S a USD\,30/ud.~a cualquier volumen representa un costo
  elevado para un dispositivo de consumo. La integración directa del chip
  WiFi (e.g., Realtek RTL8822CS o similar) en PCB eliminaría el premium del
  módulo, reduciendo este subsistema a USD\,3--5/ud., a costa de mayor
  complejidad de diseño de RF y antena. Esta decisión no afecta al E3 pero
  debe documentarse como riesgo estratégico para revisiones futuras.
\end{decisbox}

\subsection{Sensibilidad ante Variaciones de Precio}

Una fluctuación del $\pm$10\,\% en el precio del SoC (USD\,$\pm$28.57/ud.)
tiene un impacto de $\pm$0.36\,\% en el precio de venta objetivo, lo cual es
manejable. En cambio, una escasez del módulo WiFi AP6275S o del DRAM podría
incrementar el costo de materiales en USD\,5--15/ud., afectando directamente
el margen del producto. El plan de mitigación documentado en el
Capítulo~\ref{chap:riesgos} contempla la identificación de fuentes
alternativas para estos dos componentes antes del inicio del E4.

